\documentclass[11pt]{article}
\usepackage[margin=1in]{geometry}
\usepackage{graphicx}
\usepackage{booktabs}
\usepackage{amsmath}
\usepackage{amssymb}
\usepackage{natbib}
\usepackage[hidelinks]{hyperref}
\usepackage{xcolor}
\graphicspath{{../figures/}}

\title{\textbf{When Does the LLM Disambiguation Stage Actually Help?}\\
\large A decomposition of retriever--LLM interaction in biomedical entity linking}
\author{Moritz \and Mohamed \and Snehpreet\\
\small NLP Lab, University of Bonn \quad (Supervisor: Frederik)}
\date{\today}

\begin{document}
\maketitle

\begin{abstract}
Large language models (LLMs) are increasingly bolted onto biomedical entity-linking
(EL) pipelines as a final disambiguation stage, and recent work reports large gains
from doing so. We do not propose a new method; we ask \emph{how much}, \emph{when}, and
\emph{why} this stage actually helps. Using a reproduced neuro-symbolic pipeline as a
measuring instrument across three MeSH datasets (BC5CDR, BioRED, NCBI-Disease),
NLM-Chem, and MedMentions, and varying retriever strength on the same mentions, we find:
(i) the LLM's value is set by the retriever, not by the LLM---its contribution swings
roughly ten-fold with the retriever alone, and improving the retriever yields
diminishing returns because the two are partial substitutes; (ii) the zero-shot LLM is
harmful in a predictable mid-confidence ``danger zone'', which a confidence gate removes
(on BC5CDR the gate beats running the LLM everywhere while using under one-fifth of the
calls); and (iii)---our central claim---LoRA fine-tuning does not make the model a better
disambiguator where disambiguation is hard; it installs a confidence-dependent
\emph{restraint} (deference) whose reliability benefit is bound to the training
distribution, near-perfect on seen concepts and absent on unseen ones. We frame (ii) as
the entity-linking instantiation of the Confidence Gate Theorem and, in (iii), report a
phenomenon outside that theorem's static setup: fine-tuning reshapes the intervener's
own inversion zone, but only in-distribution.
\end{abstract}

\section{Introduction}
Biomedical entity linking maps a mention in text to a concept in a knowledge base such
as MeSH or UMLS. Most systems are two-stage: a retriever generates candidate concepts,
and a disambiguation stage selects among them. LLMs are now routinely used for the second
stage, and \citet{ye2025} report new state of the art (up to $+16$ points) from a
training-free LLM disambiguator, while \citet{llm4bioel2025} obtain competitive results
with constrained decoding. The field treats this stage as a semantic reasoner. Yet its
contribution is rarely isolated: a single accuracy number hides where the LLM helps,
where it hurts, and what fine-tuning changes.

We reframe the problem from ``build a better system'' to ``analyse and explain a stage
that the field has not yet characterised.'' Our contributions are:
\begin{enumerate}
\item A frame: the LLM's marginal value is a decreasing function of retriever strength,
measured as a continuous relationship across five datasets and, causally, within
datasets (Section~\ref{sec:frame}). Retriever and LLM are partial substitutes.
\item A mid-confidence \emph{danger zone}: the zero-shot LLM over-rides a
mostly-correct retriever and is net-harmful in a specific confidence band, which
replicates across four datasets; a dev-tuned confidence gate recovers this and, on
BC5CDR, exceeds full-LLM accuracy at under $20\%$ of the calls (Section~\ref{sec:f1}).
\item \textbf{Deference, not competence}: fine-tuning leaves competence flat on the hard
cases and instead learns a distribution-bound restraint---its precision benefit is
realised almost entirely on concepts seen in training (Section~\ref{sec:f2}).
\item An anatomy of the interventions: only a minority of the LLM's prediction changes
are genuine fixes, and it recovers only a fraction of the errors it could
(Section~\ref{sec:f3}).
\end{enumerate}

\section{Related Work}
\textbf{Biomedical EL and the LLM stage.} BioLinkerAI \citep{biolinkerai2024} is a
neuro-symbolic pipeline (linguistic rules, candidate generation, domain rules, LLM
disambiguation) that we reproduce as our instrument. \citet{ye2025} isolate the LLM as
a disambiguator on top of SapBERT and ArboEL and observe---qualitatively, over two fixed
base models---that its value depends on the base model's strength; we make this causal,
quantify it, and extend it to fine-tuning. \citet{llm4bioel2025} guide the LLM with
restrictive and contrastive decoding; their own ablations show the LLM machinery adds
$\leq 0.3$ points on BC5CDR, indirect support for our frame. The current state of the
art is a trained generative model, ANGEL \citep{angel2024}, at $94.7$ on BC5CDR; it has
no separable retriever stage and is orthogonal to our analysis.

\textbf{Confidence gating.} \citet{doku2026} prove that confidence-based abstention
improves a ranked system monotonically iff the confidence signal has no ``inversion
zones''. Our retriever-gap signal contains one; we give the entity-linking instantiation
and show fine-tuning reshapes it.

\textbf{Memorisation under fine-tuning.} That fine-tuning trades cross-domain
generalisation for in-domain performance is known for EL retrievers
\citep{logeswaran2019} and entity matching. Our contribution is the mechanism for the
LLM disambiguation stage, with a zero-shot control.

\section{Experimental Setup}
We use the reproduced BioLinkerAI pipeline: retrieval (Phase~2), rule-based re-ranking
(Phase~3), and LLM disambiguation (Phase~4). The disambiguator is Qwen3-4B, both
zero-shot and LoRA fine-tuned on BC5CDR. We evaluate on BC5CDR, BioRED and NCBI-Disease
(MeSH diseases/chemicals), NLM-Chem (MeSH chemicals, full text), and MedMentions (full
UMLS). We vary retriever strength on the \emph{same} mentions (lexical vs.\ dense
SapBERT retrieval). We report, per confidence band (the Phase-3 top1--top2 score gap):
the intervention \emph{rate}, intervention \emph{precision} (fixes$/$(fixes$+$breaks)),
and net Accuracy@1 change $\Delta = \text{P4}-\text{P3}$, with document-level bootstrap
$95\%$ CIs. ``Seen'' concepts are those in the fine-tuning training set.

\section{The Frame: the retriever sets the LLM's value}\label{sec:frame}
\begin{figure}[t]\centering
\includegraphics[width=.72\linewidth]{fig1_curve.pdf}
\caption{Net $\Delta$ Accuracy@1 of the LLM stage (P4$-$P3), per retriever-confidence
quintile (Q1 least $\to$ Q5 most confident; equal-count bins of the Phase-3 top1$-$top2
gap, computed \emph{per run} so columns are comparable across retrievers with different
score scales). Rows are eight zero-shot-Qwen3-4B runs sorted by Candidate Recall@1; the
right column is the overall gain. Blue $=$ the LLM helps, red $=$ it hurts. The LLM is
uniformly valuable where the retriever is unsure (Q1); a mid-confidence \emph{net-harmful}
band (red) emerges as the retriever strengthens (downward), and the overall gain shrinks
from $+7.1$ to $+0.6$. BM25 and BioSyn are standalone public retrievers we did not
build.}\label{fig:curve}
\end{figure}
The same LLM with the same prompt contributes wildly different amounts depending only on
the retriever: in the overall-gain column of Figure~\ref{fig:curve} the gain sits at $+4$
to $+7$ points for the weaker runs but collapses to $+0.9$ and $+0.6$ on the two strongest
retrievers (our full pipeline and BioSyn). This is causal, not correlational: on the
\emph{same} MedMentions mentions, strengthening only the retriever (Recall@1 $47.8\% \to
55.3\%$) halves the LLM's gain ($+8.5 \to +4.3$), and the same within-dataset pattern---a
stronger retriever, a smaller gain---holds on four of the five datasets; it is absent only
on NCBI, where dense retrieval did not improve Recall@1. The
relationship is known from IR
reranking and, qualitatively, from \citet{ye2025}; our contribution is to measure it as
a continuous curve, and it survives normalisation by the available headroom (the LLM's
recovery rate falls too, $59\% \to 30\%$ on BC5CDR). A direct consequence is that
retriever and LLM are partial \emph{substitutes}: on MedMentions a $+7.5$-point Recall@1
improvement raises final accuracy by only $+3.2$, because the LLM was already recovering
most of those cases (Figure~\ref{fig:subst}). The LLM's value is thus not intrinsic---it
is borrowed from the retriever's weakness.
\begin{figure}[t]\centering
\includegraphics[width=.7\linewidth]{fig2_substitution.pdf}
\caption{Final accuracy $=$ retriever Recall@1 (blue) $+$ LLM gain (red). As the
retriever strengthens the LLM's slice shrinks; the final barely moves.}\label{fig:subst}
\end{figure}

\section{Finding 1: a mid-confidence danger zone, and a gate}\label{sec:f1}
\begin{figure}[t]\centering
\includegraphics[width=\linewidth]{f1a_dangerzone.pdf}
\caption{Intervention rate (line) falls monotonically with retriever confidence, but
precision (bars, $95\%$ CI) collapses in the mid-high band [10,20), where the LLM is
net-harmful (red) on the two MeSH disease datasets and absent-as-predicted on
weak-retriever MedMentions.}\label{fig:danger}
\end{figure}
The LLM's intervention \emph{rate} falls monotonically as the retriever grows more
confident (BC5CDR $54.3\% \to 0.9\%$; Figure~\ref{fig:danger}, line): it sensibly acts
less when the retriever is sure. Its \emph{precision}, however, collapses in a mid-high
band rather than at the extreme: at gap [10,20) it still intervenes on $6.1\%$ of
mentions and is right only $18\%$ of the time ($95\%$ CI $[10,30]$; 20 fixes vs.\ 89
breaks), producing a net $-2.4$ points there, because the retriever is already correct on
$86.5\%$ of those mentions. The naive reading ``more confident retriever $\Rightarrow$
more harmful LLM'' is wrong: the harm is localised to a band where the model is confident
enough to act but the retriever is already reliable. This net-harmful band replicates on
BC5CDR ($-2.4$), BioRED ($-1.4$), NCBI ($-7.6$) and NLM-Chem ($-1.0$), and vanishes on
MedMentions---exactly where the mechanism predicts, since a weak retriever is not
reliable enough for the LLM to harm. In the terms of \citet{doku2026}, this is an
inversion zone in the retriever-gap confidence signal.

\begin{figure}[t]\centering
\includegraphics[width=.62\linewidth]{f1b_gate_pareto.pdf}
\caption{Confidence gate: fraction of the full-LLM gain captured vs.\ fraction of
mentions sent to the LLM ($\blacklozenge$ = threshold tuned on held-out development
documents). BC5CDR exceeds $100\%$: the gate beats running the LLM everywhere.}
\label{fig:gate}
\end{figure}
Because the harm is confined to a predictable band, a confidence gate---call the LLM only
when the gap is small---recovers it. Tuning the threshold on held-out development
documents and reporting on the rest, on BC5CDR the gate sends only $19\%$ of mentions to
the LLM yet reaches $85.5\%$ Accuracy@1, above the $85.2\%$ of running it everywhere
(Figure~\ref{fig:gate}). The gate's efficiency scales with retriever strength: under
$30\%$ of calls retains $\geq 80\%$ of the benefit on the strong-retriever MeSH datasets,
whereas MedMentions needs the LLM on roughly $60\%$ of mentions. On the fine-tuned model
the gate no longer beats full accuracy, because fine-tuning removes the harmful band
(Section~\ref{sec:f2}); ``less is more'' holds for the \emph{un-tuned} LLM.

To check these effects are not artefacts of our hand-built pipeline, we repeat the BC5CDR
analysis with candidates and confidence drawn purely from two public retrievers we did
not build: a generic SapBERT bi-encoder (Recall@1 $75.9\%$; the retriever used by
\citet{ye2025} and \citet{llm4bioel2025}) and BioSyn \citep{biosyn2020} (Recall@1
$84.8\%$), a BC5CDR-tuned dense+sparse system. On both, the value--confidence structure
replicates (rate monotone, value concentrated in low confidence, precision lowest in a
mid band), and both fall on the downward trend of Figure~\ref{fig:curve}. What differs is
exactly what the retriever-strength mechanism predicts: on the \emph{weaker} SapBERT every
band is net-positive (the gate only matches the full LLM), whereas on the \emph{stronger}
BioSyn the full danger zone reappears with a public signal---precision falls to $34\%$ and
$23\%$ in the $[5,10)$ and $[10,20)$ bands, net $-2.5$ and $-1.7$ points, while
low-confidence bands stay strongly positive ($+17.7$, $+7.2$). The gate consequently
\emph{beats} the full LLM on BioSyn ($87.0\%$ at $15\%$ of calls vs.\ $85.4\%$ full). The
self-harm and ``gate beats full'' are thus not artefacts of our rule-based score: they
appear whenever the retriever is strong enough that its mid-confidence top-1 is usually
right, and vanish otherwise. The danger zone scales with retriever strength---tying
Finding 1 to the Frame and matching the inversion-zone condition of \citet{doku2026}.

\section{Finding 2: fine-tuning buys deference, not competence}\label{sec:f2}
\begin{figure}[t]\centering
\includegraphics[width=.6\linewidth]{fig4_memorization.pdf}
\caption{Intervention precision, seen vs.\ unseen concepts, zero-shot vs.\ fine-tuned
(illustrative single-corpus view). Fine-tuning raises seen precision and does not on
unseen.}\label{fig:mem}
\end{figure}
Read by confidence band, fine-tuning does not improve the model where disambiguation is
hard: on the truly-uncertain band (gap [0,2)) precision is unchanged ($98\% \to 98\%$ on
BC5CDR, $93\% \to 94\%$ on BioRED). It instead learns to stop over-riding a confident
retriever---on BC5CDR the net-harmful band [10,20) flips from precision $18\%$ to $64\%$
and from $\Delta -2.4$ to $+0.7$, while the intervention rate there falls
($6.1\% \to 4.4\%$). This is the gate of Finding~1 learned internally, and consistently
fine-tuning halves the breaks share ($21\% \to 8\%$; Section~\ref{sec:f3}).

Whether this restraint \emph{transfers} depends on whether the concept was in the
fine-tuning distribution. Pooling three held-out corpora the model never trained on
(BioRED, NCBI, NLM-Chem; $4{,}024$ mentions), fine-tuning raises intervention precision
on seen concepts from $69\%$ to $97\%$ ($95\%$ CIs $[49,83] \to [92,99]$; non-overlapping;
breaks collapse $42 \to 5$), while on unseen concepts precision is essentially unchanged,
$72\% \to 74\%$ (CIs $[62,81] \to [65,83]$; overlapping; breaks rise $77 \to 88$). The
reliability fine-tuning buys is realised almost entirely on seen concepts; the seen-vs.-
unseen difference in the mid/high bands is large ($+24$ and $+53$ points). We make the
conservative claim the data support: fine-tuning installs a confidence-dependent
restraint whose precision benefit is \emph{bound to the training distribution}. This one
cause subsumes Findings~1 and~3: fine-tuning changes \emph{when} the stage intervenes,
not \emph{how well} it decides.

\section{Finding 3: an anatomy of the interventions}\label{sec:f3}
\begin{figure}[t]\centering
\includegraphics[width=.7\linewidth]{fig5_churn.pdf}
\caption{Decomposition of the LLM's prediction changes into fixes, breaks, and lateral
moves (no accuracy change). Only a minority are genuine fixes.}\label{fig:churn}
\end{figure}
Of the zero-shot LLM's prediction changes on BC5CDR, only $33\%$ are fixes; $21\%$ are
breaks and $46\%$ are lateral (Figure~\ref{fig:churn}). Splitting the lateral moves is
essential: of the $355$, only $112$ ($32\%$) occur when the gold concept is in the
candidate list---genuine confusion---while the remaining $243$ are forced, because the
gold is not retrievable and any change is lateral by construction. The genuine confusion
rate is therefore $\sim 15\%$ of all changes, not $46\%$. Fine-tuning raises fixes to
$51\%$ and cuts breaks to $8\%$ while the genuine confusion barely moves ($15\% \to
10\%$): it removes harmful over-rides (deference) without adding disambiguation skill
(competence). The stage is best described as \emph{over-eager where there is nothing to
gain, and under-powered where it matters}: even when the gold answer is in the candidate
list, it recovers only $\sim 20$--$42\%$ of the retriever's mistakes.

\section{Discussion}
Across five datasets one account ties the findings together. The LLM disambiguation
stage is not a general semantic reasoner applied uniformly; its value is a
retriever-dependent, mechanically bounded quantity, and it is idle-to-harmful on the
confident majority. Fine-tuning does not lift the ceiling on the hard cases---it teaches
the model \emph{when} to defer, and that deference is memorised, calibrated to the
training distribution. The practical consequences are concrete: behind a strong retriever
the stage should be gated or cascaded (cheaper and at least as accurate); investment in
the retriever has diminishing final returns because the LLM masks it; and fine-tuning a
disambiguator should be expected to help on familiar concepts and not on novel ones.

\section{Limitations}
Our results use a single open 4B disambiguator and a single fine-tuning run; seed
variance and a second training source (or multi-corpus fine-tuning) would strengthen the
deference claim. The seen/unseen split is confounded with concept frequency, which a
within-frequency stratification would isolate. The gate threshold is tuned on a single
development split and should be read with the full Pareto curve. To rule out that the
findings are properties of our implementation rather than of LLM disambiguation, the
decomposition should be applied to a public system we did not build (e.g.\ LLM4BioEL),
and ``recovery is weak'' should be calibrated against a cross-encoder baseline on the
same candidate lists. Our retriever axis pools configurations and datasets; a denser
sweep with named public retrievers (BM25, SapBERT, BioSyn, ArboEL) is in progress.

\section{Conclusion}
Isolating the LLM disambiguation stage in biomedical EL shows that its value is set by
the retriever, that it is harmful in a predictable mid-confidence band a gate removes,
and that fine-tuning buys deference rather than competence---a confidence-dependent
restraint bound to the training distribution. These are properties the field's single
accuracy numbers hide, and they carry direct guidance for when to deploy, gate, or
fine-tune the stage.

\begin{thebibliography}{9}
\bibitem[BioLinkerAI(2024)]{biolinkerai2024} BioLinkerAI. \emph{A neuro-symbolic approach
for biomedical entity linking.} Springer, 2024.
\bibitem[Ye \& Mitchell(2025)]{ye2025} C.\ Ye and C.\ S.\ Mitchell. \emph{LLM as Entity
Disambiguator for Biomedical Entity-Linking.} ACL (Short), 2025.
\bibitem[Lin et al.(2025)]{llm4bioel2025} Z.\ Lin, Z.\ Zhang, J.\ Wu, Y.\ Zheng, X.\ Wu.
\emph{Guiding LLMs for Biomedical Entity Linking via Restrictive and Contrastive
Decoding.} Findings of EMNLP, 2025.
\bibitem[Kim et al.(2024)]{angel2024} Kim et al. \emph{Learning from Negative Samples in
Biomedical Generative Entity Linking (ANGEL).} arXiv:2408.16493, 2024.
\bibitem[Doku(2026)]{doku2026} R.\ Doku. \emph{The Confidence Gate Theorem: When Should
Ranked Decision Systems Abstain?} arXiv:2603.09947, 2026.
\bibitem[Logeswaran et al.(2019)]{logeswaran2019} L.\ Logeswaran et al. \emph{Zero-Shot
Entity Linking by Reading Entity Descriptions.} ACL, 2019.
\bibitem[Sung et al.(2020)]{biosyn2020} M.\ Sung, H.\ Jeon, J.\ Lee, and J.\ Kang.
\emph{Biomedical Entity Representations with Synonym Marginalization.} ACL, 2020.
\end{thebibliography}

\end{document}
